Der Ofen muss gerade so viel Leistung abgeben, wie durch die Wände entweicht. Für eine Massivholzwand der Dicke \dHO\ gibt FoTaBe (S.~190) den Wärmedurchgangskoeffizienten $U = \UHO$. Mit $\Delta T = \dTF = \dT$ ist
\begin{align}
 P &= \PHF \\
   &= \UH\cdot\AH\cdot\dT \\
   &= \PH \approx \result{\PHP{3}{2}}
\end{align}
Rechnet man nur mit der Wärmeleitfähigkeit von Holz ($\lambda = \lHO$ für Tannentäfer, FoTaBe S.~190), erhält man $P = \PLF = \PLP{3}{2}$. Dieser Wert ist zu gross, weil er die Wärmeübergangswiderstände an den beiden Wandoberflächen nicht berücksichtigt. Verluste durch Fenster, Türen, Dach, Boden und Lüftung sind in beiden Rechnungen nicht enthalten.
